\documentclass[10pt]{article}
\usepackage{icmlko}

\icmlmeta{IP 파운데이션 월드모델}{276}{문을 열었더니 안에 문이 하나 더}{팝업 전용 축 넷이 판·도메인 전부 정확히 0.0000 이었다 --- 자료가 아니라 나무의 잎 하한이 막고 있었다}

\begin{document}

\twocolumn[
\icmlheader

\begin{abstract}
노트 275가 학습 문턱을 40에서 15로 내려 팝업을 학습 풀에 넣었다. 그래서
노트 276은 팝업 \textbf{전용} 축을 만들었다 --- 프로젝트 레코드 380건에서
후보 열여섯을 걸러 넷을 남겼다(경쟁 열 넷이 서로 0.85$\sim$1.00 이라
\texttt{500m} 과 \texttt{1km} 은 순위 상관이 \textbf{정확히 1.000}, 하나만
썼다. \texttt{기간\_일수}는 검사를 통과했는데도 뺐다 --- 라벨 \texttt{y\_perday}
의 \textbf{분모}였다). 재니 판도 아홉 도메인도 \textbf{전부 정확히 0.0000},
짝 SE 까지 0 이었다. 노트 187의 제일 값싼 검사가 걸린 것이다 --- \emph{효과가
없는 게 아니라 모형에 안 닿았다.} 라벨 순위를 그대로 전용 열로 줘 봐도
\textbf{0.0000} 이었다. \textbf{그런데 배선은 살아 있었다} --- 설계 자리
23$\sim$26에 있고 팝업 학습 16행 중 14행에서 값이 변한다. 도메인마다 같은
누출 열을 줘 보니 학습행 54(아이돌)까지는 $+$0.35$\sim$0.70 을 얻고 16(팝업)
에서만 0 이다. 벽이 그 사이에 있다. \textbf{범인은 \texttt{min\_samples\_leaf}
$=$ 20}(sklearn 기본값, 우리가 한 번도 안 정한 값)이다. 관측 14행으로는 어느
쪽 잎도 20을 못 채우니 그 열로 가르는 분기가 \textbf{원리적으로 무효}다.
하한만 20 $\to$ 10 $\to$ 5 로 내리니 같은 누출 열이 \textbf{$+$0.0000 $\to$
$+$0.0394 $\to$ $+$0.4373} 으로 살아난다. \textbf{벽은 자료가 아니라
손잡이였다.} 그리고 \texttt{MIN\_TRAIN}$=$40 과 \texttt{min\_samples\_leaf}
$=$20 은 사실 같은 상수다 --- 양쪽 잎을 채우려면 40이 필요하다. 노트 275가
연 문 안에 문이 하나 더 있었고, 아무도 그 문을 정한 적이 없다.
\end{abstract}
]

\section{전용 축 넷을 만들었다}

노트 274가 팝업에 대해 둘을 확정했다 --- 판은 팝업을 못 본다(도메인 안에서
$\Delta\rho{=}0.435$ 가 있어야 한다), 그래도 배포 이유는 있다. 노트 275가
그걸 규칙으로 묶고 문턱을 15로 내려 팝업을 학습 풀에 넣었다.

그래서 팝업 프로젝트 레코드 380건에서 아직 축이 아닌 파생 열을 봤다. 쓰는
행 75건 위에서 덮음 88$\sim$91\% 인 후보가 열여섯이었고, 노트 239 · 256의
채택 검사로 걸렀다. \textbf{둘이 잡혔다.}

\begin{center}\small
\begin{tabular}{ll}
\hline
걸린 것 & 왜 \\
\hline
경쟁 열 넷 & 서로 0.85$\sim$1.00. \texttt{500m}$\times$\texttt{1km} $=$
\textbf{1.000} \\
달력\_연휴전 & \texttt{cal\_holiday\_gap} 과 0.592 --- 다시 쓰기 \\
\textbf{기간\_일수} & 검사 통과(0.177)했지만 라벨의 \textbf{분모} \\
\hline
\end{tabular}
\end{center}

\texttt{기간\_일수}는 \texttt{state/dataset\_v2.py} 가 이미 답을 적어 뒀다 ---
``\texttt{planned\_operating\_days} 는 피처가 아니라 \textbf{분모}로
승격시킨다; \texttt{y\_perday} $=$ visitors / days.'' \textbf{통과가 근거는
아니다}(노트 133). 남긴 넷: \texttt{pop\_comp} · \texttt{pop\_wkend} ·
\texttt{pop\_holid} · \texttt{pop\_stores}.

\section{전부 정확히 0.0000}

\begin{figure*}[t]
\centering
\includegraphics[width=\textwidth]{figs/floor.pdf}
\caption{왼쪽 --- 도메인마다 \textbf{라벨 순위를 그대로} 전용 열로 주고 그
도메인 $\rho$ 가 얼마나 오르는지 잰 것. 학습행 54까지는 크게 오르고 16에서만
0 이다. 붉은 띠가 \texttt{min\_samples\_leaf}$=$20 아래. 오른쪽 --- 팝업에
같은 누출 열을 주고 잎 하한만 바꾼 것.}
\end{figure*}

\begin{center}\small
\begin{tabular}{lrrr}
\hline
& 23축 & 27축 & 차 \\
\hline
판 (문턱 40) & 0.4845 & 0.4845 & $+$0.0000 \\
판 (문턱 15) & 0.4857 & 0.4857 & $+$0.0000 \\
팝업 (문턱 15) & 0.3931 & 0.3931 & $+$0.0000 \\
\hline
\end{tabular}
\end{center}

아홉 도메인 전부 넷째 자리까지 같고 \textbf{짝 SE 도 0.0000} 이다 --- 예측이
바이트 단위로 같다는 뜻이다. 노트 187이 적은 대로, 이건 ``효과가 없다''가
아니라 \textbf{모형에 안 닿은 것}이다.

\paragraph{첫 진단이 틀렸다} 라벨 순위를 그대로 전용 열로 넣어 봤다. 이것도
\textbf{$+$0.0000} 이길래 ``배선이 죽었다''고 적었다. \textbf{틀렸다.} 그
검정은 ``안 닿는다''와 ``닿지만 못 쓴다''를 못 가른다. 설계행렬을 직접 열어
보니 배선은 멀쩡했다.

\begin{center}\small
\begin{tabular}{lr}
\hline
확인한 것 & \\
\hline
문턱 15에서 팝업이 학습에 있나 & \textbf{있다} (16행) \\
\texttt{axis\_order} 가 골랐나 & \textbf{골랐다} (자리 23$\sim$26) \\
그 열이 학습행에서 변하나 & \textbf{변한다} (14행 · 값 14가지) \\
팝업의 학습 몫 & 16 / 10,292 $=$ \textbf{0.16\%} \\
\hline
\end{tabular}
\end{center}

\section{54행은 배우고 16행은 못 배운다}

그러면 몇 행부터 되나. \textbf{도메인마다} 자기 라벨 순위를 전용 열로 하나씩
주고 쟀다. 라벨이 새는 열이니 배울 수만 있으면 크게 올라야 한다.

\begin{center}\small
\begin{tabular}{lrrrr}
\hline
도메인 & 학습행 & 몫 & 누출 $\rho$ & 이득 \\
\hline
세계애니 & 2,648 & 25.7\% & 0.9911 & $+$0.4557 \\
웹툰 & 2,106 & 20.5\% & 0.9938 & $+$0.5490 \\
모바일 & 1,559 & 15.2\% & 0.9873 & $+$0.4444 \\
애니 & 1,467 & 14.3\% & 0.9818 & $+$0.4810 \\
펀딩 & 320 & 3.1\% & 0.9581 & $+$0.6992 \\
게임 & 259 & 2.5\% & 0.9755 & $+$0.3561 \\
도서 & 80 & 0.8\% & 0.7400 & $+$0.3472 \\
아이돌 & 54 & 0.5\% & 0.7063 & $+$0.5677 \\
\textbf{팝업} & \textbf{16} & \textbf{0.16\%} & 0.3819 & \textbf{$+$0.0000} \\
\hline
\end{tabular}
\end{center}

\textbf{비탈이 아니라 벼랑이다.} 54행까지는 다 배우고(누출 $\rho$ 0.71 이상)
16행에서 \emph{정확히} 0 이다. 몫이 0.5\% 여도 되는데 0.16\% 는 안 된다 ---
몫이 원인이면 이렇게 안 끊긴다.

\section{범인은 우리가 안 정한 값이었다}

\texttt{F18\_bagboost} 는 \texttt{HistGradientBoostingRegressor} 를 쓰는데
설정이 이렇다.

\begin{center}\small
\begin{verbatim}
self.kw = dict(max_depth=depth, max_iter=it,
               learning_rate=lr,
               l2_regularization=1.0, random_state=0)
\end{verbatim}
\end{center}

\texttt{min\_samples\_leaf} 가 없다. \textbf{sklearn 기본값 20}이다. 팝업의
관측 학습행은 \textbf{14}다. 그 열로 어디를 가르든 한쪽 잎이 14 이하이므로
\textbf{그 분기는 후보에도 못 오른다.} 아이돌 54 · 도서 80 은 20을 넘는다.
\textbf{벼랑이 54와 16 사이에 있는 이유가 이것이다.}

하한만 바꿔 확인했다.

\begin{center}\small
\begin{tabular}{rrrr}
\hline
잎 하한 & 팝업 정상 & 팝업 누출 & 이득 \\
\hline
20 (기본값) & 0.3819 & 0.3819 & \textbf{$+$0.0000} \\
10 & 0.3990 & 0.4383 & $+$0.0394 \\
5 & 0.3861 & 0.8234 & \textbf{$+$0.4373} \\
2 & 0.3832 & 0.8618 & $+$0.4786 \\
\hline
\end{tabular}
\end{center}

\textbf{하한을 5로 내리면 같은 열이 $+$0.4373 을 낸다} --- 큰 도메인들이
받는 크기와 같다. \textbf{자료가 모자란 게 아니라 손잡이가 막고 있었다.}

\paragraph{두 상수가 사실 하나였다} \texttt{MIN\_TRAIN}$=$40 은 노트 5의
편의값이라고 적혀 있고 \texttt{min\_samples\_leaf}$=$20 은 우리가 정한 적이
없다. 그런데 \textbf{40 $=$ 2$\times$20} 이다 --- 한 도메인이 제 열로 가르려면
양쪽 잎을 채워야 하니 최소 40행이 필요하다. \textbf{문턱 40은 우연히 옳았고,
노트 275가 그걸 15로 내렸을 때 안쪽 문은 그대로였다.} 팝업은 방에 들어왔지만
제 열로는 말을 못 한다.

\section{그래서 무엇을 하나}

\textbf{하한을 내리지는 않는다.} 위 표에서 판이 0.4842 $\to$ 0.4824 $\to$
0.4839 $\to$ 0.4828 로 내리거나 제자리다 --- 팝업 하나 살리자고 아홉을
흔드는 것은 노트 275가 세운 규칙의 ②(어느 도메인도 $|t|{>}2$ 로 안 나빠질 것)
를 재기 전에는 못 한다. 그리고 이건 원래 \textbf{과적합을 막으라고} 있는
값이다.

\textbf{대신 규약을 고친다.} 노트 250 · 251이 ``전용 이름으로 준다''를 세울
때 근거는 웹툰(2,106행)이었다. \textbf{그 규약은 큰 도메인에서만 맞다.}

\begin{quote}\small
학습행이 \texttt{2$\times$min\_samples\_leaf} 아래인 도메인에는 전용 이름을
주지 않는다. 그 도메인의 레코드 필드는 \textbf{공유 축 이름에 태워야} 한다.
\end{quote}

근거는 노트 275에 이미 있었다 --- 문턱을 내려 팝업이 얻은 $+$0.0317 은
\emph{공유} 열의 적합이 팝업 16행을 흡수해서 난 것이다. \textbf{팝업의 행은
공유 열로만 일한다.} 오늘 만든 전용 축 넷은 그래서 0 이었다.

\paragraph{남은 것}
\begin{center}\small
\begin{tabular}{ll}
\hline
갈래 & 왜 \\
\hline
넷을 공유 축에 태우기 & \texttt{pop\_comp} 를 \texttt{venue\_prominence} 로 \\
잎 하한을 손잡이로 & 노트 275 규칙 ②로 정식 판정 \\
아이돌 54행 & 벼랑 바로 위. 전용 축이 실제로 될까 \\
다른 정식화 & F21 능형에는 잎 하한이 없다 --- 거기선? \\
\hline
\end{tabular}
\end{center}

\paragraph{규약} \textbf{정확히 0.0000 이 나오면 자료를 의심하기 전에 기본값을
읽는다.} 노트 187이 ``안 움직이면 안 닿은 것''을 세웠는데, 오늘은 닿았는데도
안 움직였다 --- 한 겹 더 있다: \textbf{닿았지만 쓸 수 없는 경우}. 그리고
\textbf{우리가 정한 적 없는 값이 우리 결론을 정하고 있는지 본다} --- 이번에
그 값은 라이브러리 기본값이었고, 이 실험실이 세 노트에 걸쳐 ``출발점은
구조적으로 못 배운다''고 쓸 뻔했다.

\end{document}
